\chapter[Sermon 95. Furthermore Yet Is Described the Everlasting Country in Heaven.]{Furthermore Yet Is Described the Everlasting Country in Heaven.}
\label{sermon:095}
\markright{Sermon 95. Furthermore Yet Is Described the Everlasting Country ...}

And I saw no temple in it. For the Lord God Almighty and the Lamb are its temple, and the city has no need of the sun, neither of the moon, to light it. For the brightness of God illuminates it, and the Lamb is its light. And those who are saved shall walk in its light, and the kings of the earth shall bring their glory and honor into it. And its gates shall not be shut by day, for there shall be no night there. And they shall bring the glory and honor of the Gentiles into it. And there shall enter into it none unclean, neither whatever works abomination, or makes lies, but those who are written in the Lamb’s book of life.

\index{Apocalypse (Book of)}\index{John, Saint}The Apocalypse proceeds in the description of the divine or celestial City, to comfort and keep the faithful in all temptations and afflictions. Therefore, in the seventh place, he discusses the temple. For in famous cities, there is no small consideration and praise of churches. This is evident by all writers of histories, places, and times. What temple then is in heaven? None at all. For Saint John says, “And I saw in the city of God no temple.” This passage does not contradict those things which are in the 11th and 15th chapters concerning the temple in heaven. For the temple is there exhibited in a figure and vision, not that there is in deed any temple in heaven, but that thus God’s justice and certain salvation promised in the Scriptures might be signified, as we have declared in those places.

And what is the cause that no temple appears in heaven? The divine revelation answers: for the Lord God Almighty, and the Lamb are the temple in that our heavenly country. The purpose of temples is this: the Lord first instituting the tabernacle, and then the temple, wanted it to be known that he would be present in the midst of his people, as a father, Lord, and defender. And therefore it is said in the scriptures that people come unto the Lord, whether they came to the tabernacle or the temple of the Lord. The temple moreover was erected for preaching and prayer and the external service of God, for receiving the sacraments, or offering up sacrifices. But the Saints in the heavenly country have no need of all these things. Therefore they need no temple. Therefore no temple is seen in heaven. For the Lord God now shows himself to them to be enjoyed; the saints are now with him, wherefore they need no token of his presence. We are taught by doctrine what God is, and what is his will, and that we are saved by the Lamb: but now that we see God himself face to face, and that salvation has come by the Lamb of God, what need is there of a temple in heaven? By prayer we require life and everlasting joys: now since these have happened to the elect, what need is there of a house of prayer? The Saints now, without any temple, offer up eternal praises unto God. And seeing that sacrifices and sacraments have no further place in the everlasting country, I see not why there should be any temple in Heaven. We rest and keep in Heaven an everlasting Sabbath. This place moreover proves that Christ is very God, coequal with the Father, as to whom he is joined inseparably in all glory. Neither is the Holy Spirit separated from the Father and the Son, which elsewhere is said to dwell in us: for this cause we are called the temples both of God and of the Holy Spirit, by the Apostle, in 1 Corinthians 3 and 2 Corinthians 6.

\index{Augustine, Saint}The eighth section of this description is repeated concerning the heavenly light, and that not without great cause; indeed, the same light is again commended in chapter 22. For in buildings, there is nothing more excellent than light. Otherwise, without light, all things are blind. Furthermore, he does not say that the Sun and Moon should no longer exist, but that the City of God should not need those lights. He shows the reason: for the glory of God has illuminated it. And the glory of God is the divine, celestial, and unspeakable brightness of his inapproachable light, which he inhabits, and according to his good pleasure communicates to the elect. The Lord Christ (who here is called the Lamb, for the mystery of redemption) illuminates the blessed. For by him we are clarified, and enjoy that eternal, most beautiful, and celestial light. John has borrowed this passage from chapter 60 of \textit{Confessions}, where we read: “The sun shall not be there, for the light of the day. And the brightness of the moon shall not shine there: but the Lord shall be to you a perpetual light, and your God shall be your brightness.” Furthermore, the seats of the blessed are thought to be fixed above the sphere of the Sun and Moon, and also the brightness of the saints to surpass far the light of the Sun and stars. Augustine has testified this also in chapters 24 and 30. To God almighty and eternal light, be praise and thanksgiving, who has prepared such great things for us, and gives us gifts such as no tongue can express.

\index{Antichrist}He shows in places more than one who are partakers of that light, or who are citizens of this celestial city, and what is the state of the citizens. All nations and people saved are citizens of the eternal country. Here are two things to be noted. First, that the Gentiles are made inheritors of glory, and that without any choice. For here excels not the Jew, nor the Greek, neither Roman, nor Barbarian. Again, yet not all without respect, and confusedly obtain everlasting light, but the saved only: that is to say, whom Christ has saved and redeemed from sin, the Devil, Antichrist, and from the curse and the world. And Christ saves the elect and faithful. They therefore shall in deed be partakers of the light: these are the citizens of the country everlasting. But what is their state and inheritance? They shall walk in the light of God the Father and the Lamb: that is to say, they shall have the fruition of the light and of God himself, to their joyful sweetness and fulfillment. For it is a figurative speech, to walk in the light, for that which is to enjoy light. Verily in Psalm 88 is read with a figure not much unlike: Lord, they shall walk in the light of thy countenance. And again: Thou shalt make known unto me the path of life, the fulfillment of joys is in thy sight, and gladness in thy right hand forevermore.

\index{Daniel (Book of)}\index{Ezekiel (Book of)}But especially the places in heaven, and in that divine Palace, are for Kings. Kings are governors and captains of the people, as they are called Kings and Princes, governors, Magistrates, rulers, both of the political and ecclesiastical government, Doctors, masters, teachers, Artificers, and Parents. For their duty is virtuously to govern their subjects, scholars, or children, to keep them under awe or discipline, to chastise, and direct them to the duties of life and all godliness. This if they do, they shall have a worthy place prepared in Heaven. For Daniel also says in the twelfth chapter: “But the teachers shall shine as the brightness of the firmament; and they that bring many to righteousness, as the stars everlastingly.” O therefore, O happy are you if you bring many to execute the office of righteousness. But woe to you Princes, and teachers, and masters, and Parents, if herein you be negligent. There is prepared for you in hell a place most horrible and miserable, as also Ezekiel has testified.

But if Kings have their place, and the same right honorable, in Heaven, wherefore do the Anabaptists teach, nay why do they lie, that a Christian cannot execute the office of a Magistrate? For here are Kings mentioned to be in heaven, not only as men, but as they were kings, that is, as they were good kings, and executed their office duly, and not forsaking their place, have lived a private life. For it follows, they shall bring their glory and honour unto it. And what is that glory, and what is the honour? it follows again: and they shall bring the glory and honour of nations into it: that is to say, they shall bring into heaven with them, the very nations, their people and subjects, whom they have helped in true godliness and salvation, in teaching, correcting, defending, alluring or drawing, \&c. And these be their glory and honour, for Paul in the second chapter of the first to the Corinthians says, “for we are your glory, as you shall be ours also in the day of our Lord Jesus.” And again in the first to the Thessalonians, the same Apostle says: “for what is our hope, joy, or crown of rejoicing? Are not you, in the sight of our Lord Jesus Christ, at his coming? For you are our glory and joy.”

\index{Aquinas, Thomas}Aquinas rightly says: Paul speaks after the manner of conquerors, who bring their spoils into cities. Therefore he feigns that Princes, preachers, and parents bring with them into heaven such as they have won, which to them shall be an honor and glory. These things always let us think upon, and do our duty enjoined us of God, which we perceive in the everlasting country to have so great a reward. For it shall be the greatest glory that may be, to stand with so many won in the presence of the eternal God, Lamb, and all saints. Contrariwise the greatest shame to stand with so great a multitude of men lost, and that lost through our fault and negligence. Read what things are written in the first chapter of the book of Wisdom, \&c.

\index{Arethas of Caesarea}In the tenth place follows the guarding of the heavenly gates. Certainly, in great cities, great and diligent watching, guarding, and attention are taken to ensure the gates are shut and opened at the proper time. But in heaven, such carefulness is not needed. The reason is: the gates are not usually shut during the day, but at night. But in the eternal realm, there is no night, therefore the gates are never shut. There is undoubtedly no night, but continual day. There is no treachery, no ambushes or plots, no perils or dangers: all things, in general, are safe, peaceful, quiet, secure, and sure. The same things are also read in Isaiah, but in a different sense. Arethas says: here is a double understanding, for either he means that there will be peace and security, so great that it will not need to protect the city by shutting its gates. Or else that there also the godly gates of the apostolic doctrine are open to all, for their learning, and so on. Certainly, they will need no teachers nor guides, who see all mysteries clearly now and are brought into heaven itself.

And especially cleanliness in cities is highly commended, if nothing appears that offends the sight, hearing, and smelling—nothing loathsome to look upon and to be abhorred. And in private houses, the chief praise is if all things shine, and everything stands in order, and does not lie scattered and stink.

Now therefore, in the eleventh place, he shows that there will be nothing in heaven that is offensive, that is to say, which will not be pleasing and delightful, complete and absolute. This same passage must also be understood in reference to individuals. For it follows: “except those who are written in the Lamb’s book of life.” We understand therefore, that no whoremongers, idolaters, liars, deceivers, or anything unclean, and not purged by the blood of the Son of God through faith, will enter the kingdom of heaven. The Apostle affirms this in 1 Corinthians 5 and 6, and in Ephesians 5. David also asks: “Lord, who shall dwell in your tabernacle, or who shall rest in your holy hill?” And he answers immediately: “He who walks without spot, and does righteousness,” and what follows in Psalm 15. Finally, here will be fulfilled those things written in Deuteronomy 23, concerning those who are prohibited from entering the church. Therefore, this passage has a secret doctrine and a private admonishment, instructing us that if we desire to be heirs of the everlasting kingdom, we should all, while we live on Earth, apply ourselves to righteousness and innocence. For it will follow in Revelation 22: “Outside are dogs and sorcerers and whoremongers,” and so forth. May the Lord bring us by the way of righteousness to everlasting life.
